\documentclass[fc]{tarea}

\usepackage[style=mexican]{csquotes}
\usepackage{mathtools} 
\usepackage{turnstile} 
\usepackage{mathrsfs}

\newcommand{\mturnstile}[2]{\mathrel{\vDash\rotatebox[origin=c]{180}{$\vDash$}_{#1}^{#2}}}
\newcommand{\menos}{\, \backslash \,}
\newcommand{\al}{\mathscr{A}}
\newcommand{\lan}{\mathscr{L}}
\newcommand{\sis}{\text{SF}}
\newcommand{\expresiones}{\text{Exp}}
\NewDocumentCommand{\expr}{O{\al}}{\expresiones(#1)}

% Conjuntos
\providecommand\st{}
\newcommand\SetSymbol[1][]{%
  \nonscript\:#1\vert
  \allowbreak
  \nonscript\:
  \mathopen{}}
\DeclarePairedDelimiterX\set[1]{\{}{\}}{%
  \renewcommand\st{\SetSymbol[\delimsize]}#1}
\DeclarePairedDelimiterX\bset[1]{\lbrack}{\rbrack}{#1}

\examen{Tarea 2}
\prof{Karina G. Buendía y José Dosal}
\materia{Conjuntos y lógica}
\alum{}

%\xsimsetup{solution/print = true}

\begin{document}
\HojaExamen{}{e}
\chapter*{Lógica de primer orden}

\begin{exercise}
Sea $\mathscr{L}$ el lenguaje de primer orden que tiene una constante $c_0$, 
una letra relacional binaria $\mathcal{R}^2$, una letra funcional $f^2$ 
y dos símbolos funcionales $g^3$ y $h^3$. Es decir:

\[
\mathscr{L} = \set{ \mathcal{R}^2, f^2, g^3, h^3, c_0 }.
\]

La interpretación en el dominio es $\mathbb{N}$, donde la relación entre 
elementos de los naturales es el menor estricto. El primer símbolo funcional 
se interpreta como el sucesor $S : \mathbb{N} \to \mathbb{N}$ tal que 
$S(n) = n + 1$ para $n \in \mathbb{N}$. Mientras que $g^3$ y $h^3$ son la suma 
y el producto en $\mathbb{N}$. Esta estructura se denomina 

\[
\mathscr{N} = (\mathbb{N}, <, S, +, \cdot, 0).
\]

Traduce las sentencias de $\mathscr{N}$ en fórmulas del lenguaje $\mathscr{L}$.

\begin{tasks}
    \task La relación $<$ en $\mathbb{N}$ es transitiva.
    \task La relación $<$ en $\mathbb{N}$ es antisimétrica.
    \task La relación $<$ en $\mathbb{N}$ es antirreflexiva.
    \task El producto en $\mathbb{N}$ es conmutativo.
    \task El producto en $\mathbb{N}$ es asociativo.
    \task Todo número natural tiene un sucesor.
    \task Si $n < m$, entonces $n + r < m + r$, para cualesquiera $m, \, n, \, r \in \mathbb{N}$.
\end{tasks}
    
Calcula el valor de verdad de la fórmula de $\mathscr{L}$ bajo esta interpretación.

\begin{tasks}
    \task $\exists y \, \forall x \, \mathcal{R}^2(x, y)$
    \task $\forall x \, \exists y \, \mathcal{R}^2(x, y)$
    \task $\forall x \, \forall y \, (\mathcal{R}^2(x, y) \Rightarrow \mathcal{R}^2(f^2(x), f^2(y)))$
    \task $\forall x \, \forall y \, \mathcal{R}^2(x, y) \Rightarrow \forall x \, \forall y \, \mathcal{R}^2(f^2(x), f^2(y))$
\end{tasks}

Supongamos que la función asignación $s : VAR \to \mathbb{N}$ con las siguientes asignaciones: $s(x) = 5$, $s(y) = 2$ y $s(z) = 1$.
Decide si $\mathscr{N} \vDash \varphi[s(x, y, z)]$.

\begin{tasks}
    \task $\varphi = g^3(f^2(z), y) \approx x$
    \task $\varphi = \exists w \, (\mathcal{R}^2(f^2(w), x))$
    \task $h^3(z, z) \approx y$
    \task $\forall w \, h^3(w, z) \approx w$
\end{tasks}

\end{exercise}

\begin{exercise}
    Decide si $\mathscr{A} \vDash \varphi[s]$ para las siguientes fórmulas:

\begin{tasks}
    \task $\exists x \, (\mathcal{R}(x, x) \Rightarrow \neg \mathcal{R}(x, x))$
    \task $\mathcal{R}(f(c), c) \Rightarrow \forall x \, \mathcal{R}(f(x), x)$
    \task $\exists x \, \mathcal{R}(x, x) \Rightarrow \forall x \, \mathcal{R}(x, x)$
\end{tasks}
\end{exercise}

\begin{exercise}
    Sean $\mathscr{A}$ una estructura, $\varphi(x_1, \dots, x_n)$ una fórmula y $s_1, s_2 : VAR \to A$ dos asignaciones. Si $s_1(x_i) = s_2(x_i)$
    para todo $i \in \set{1, \dots, n}$, entonces $\mathscr{A} \vDash \varphi(x_1, \dots, x_n)[s_1]$ si y solo si 
    $\mathscr{A} \vDash \varphi(x_1, \dots, x_n)[s_2]$.
\end{exercise}

\begin{exercise}
    Demuestra las siguientes afirmaciones:

\begin{tasks}
    \task $\exists x \, \mathcal{R}(x) \nvDash \mathcal{R}(x)$
    \task $\mathcal{R}(x) \vDash \exists x \, \mathcal{R}(x)$
    \task $\forall y \, \exists x \, \mathcal{R}(x, y) \nvDash \exists x \, \forall y \, \mathcal{R}(x, y)$
    \task $\forall x \, \forall y \, \mathcal{R}(x, y) \vDash \forall y \, \mathcal{R}(x, y)$
    \task Si $\vDash \alpha$ entonces $\vDash \forall x \, \alpha$
    \task $\exists x \, \mathcal{R}(x)\; \mturnstile{}{}\; \neg \forall x \, \neg \mathcal{R}(x)$
    \task $\forall x \, \mathcal{R}(x)\; \mturnstile{}{}\; \neg \exists x \, \neg \mathcal{R}(x)$
\end{tasks}

\end{exercise}

\chapter*{Teoría de conjuntos}

\begin{exercise}
    Para todo conjunto $A$, $B$ y $C$ se cumple lo siguiente:
    \begin{tasks}
    \task $A\, \triangle\, \varnothing = A$
    \task $A\, \triangle\, A = \varnothing$
    \task $A\, \triangle\, B = B\, \triangle\, A$
    \task $(A\, \triangle\, B)\, \triangle\, C = A\, \triangle\, (B\, \triangle\, C)$
    \task $A\, \cap\, (B\, \triangle\, C) = (A\, \cap\, B)\, \triangle\, (A\, \cap\, C) $
    \task Si $A\, \triangle\, B = A\, \triangle\, C$, entonces $B = C$
    \end{tasks}
\end{exercise}

\begin{exercise}
    \begin{tasks}
        \task Demuestre que si $A \subseteq C$, entonces $A \cup (B \cap C) = (A \cup B) \cap C$ .
        \task ¿Será cierto el resultado anterior si se suprime la hipótesis $A \subseteq C$?
        \task Demuestre que $A \subseteq C$ si y solo si $A \cup (B \cap C) = (A \cup B) \cap C$.
    \end{tasks}
\end{exercise}

\begin{solution}
    Para la demostras la idea tenemos que demostral la doble contención. Primero demostramos que $A \cup (B \cap C) \subseteq (A \cup B) \cap C$.
    Sea $x \in A \cup (B \cap C)$, entonces $x \in A$ o $x \in (B \cap C)$. Caso (1) supongamos que $x \in A$ y por hipótesis
    $A \subseteq C$, por lo que $x \in C$. Ademmás, como $x \in A$, entonces $X \in A \cup B$. Por lo tanto $ x \in (A \cup B)
    \cap C$. La otra contención se demuestra de manera similar.
    Para el regreso. Sea $x \in A$, entonces $x \in A \cup (B \cap C) = (A \cup B) \cap C$. Así, $x \in (A \cup B)$
    y $x \in C$.
\end{solution}

\begin{exercise}
    Sea $A$ un conjunto. Demuestre que el ``complemento'' de $A$ no es un conjunto. (El ``complemento'' de $A$ es el conjunto de todos los $x \not\in A$).
\end{exercise}

\begin{exercise}
    Pruebe lo siguiente:
\begin{tasks}
    \task $A \menos B = (A \cup B) \menos B$
    \task $A \menos (B \menos C) = (A \menos B) \cup (A \cap C)$
    \task $(A \menos C) \menos (B \menos C) = (A \menos B) \menos C$
    \task $(A \menos C) \cup (B \menos C) = (A \cup B) \menos C$
    \task $(A \menos C) \cap (B \menos C) = (A \cap B) \menos C$
    \task $(A \menos B) \menos (A \menos C) = A \cap (C \menos B)$
\end{tasks}
\end{exercise}

\begin{exercise}
    Demuestra las siguientes propiedades:
    \begin{tasks}
    \task $A_1 \cup A_2 \cup \dots \cup A_n = (A_1 \menos A_2) \cup \dots \cup (A_{n-1} \menos A_n) \cup (A_n \menos A_1) \cup \Bigl( \bigcap_{k=1}^n A_k \Bigr)$
    \task Si $A, B \subseteq X$, entonces $(X \menos A) \menos (X \menos B) = B \menos A$
\end{tasks}
\end{exercise}

\end{document}